\section{Discussion and Managerial Insights}
\label{sec:discussion}

This section interprets the empirical findings through the dual lenses of algorithmic properties and commercial fleet economics, contextualizing computational trade-offs, analytical properties, and limitations of learning-augmented metaheuristics.

\subsection{Scale-Aware Performance Dynamics: Quality-Maximization vs. Search Speed}
\label{sec:disc_scale_aware}

A central empirical finding emerging from our evaluation is that the performance advantage of Tri-Level Hybrid DDQN-ALNS exhibits a scale-aware divergence:
\begin{itemize}
    \item \textbf{100- and 200-Customer Scale (NV-Floor Convergence \& Distance Dominance)}: On Solomon-100 and Homberger-200, both classical ALNS and Hybrid-DDQN reliably converge to the minimum vehicle count floor (e.g., $NV=10.0$ on C1, $NV=6.0$ on C2\_200). In this regime, the hybrid advantage manifests as solution consistency and travel distance minimization: Hybrid-DDQN achieves travel distance savings across benchmark families, reducing the overall distance gap to BKS to $+1.41\%$ versus $+3.73\%$ for ALNS-Base. However, on complex mixed-distribution instances such as $rc1\_2\_5$, Hybrid-DDQN exhibits a noticeable $+17.11\%$ distance gap to BKS despite successfully matching the BKS vehicle count floor ($NV=18.0$ vs.\ $NV=19.0$ for ALNS-Base). This distance inflation arises because strictly eliminating the 19th vehicle in mixed topologies forces the consolidation of dispersed customer clusters under tight delivery deadlines, resulting in elongated route detours that standard neighborhood operators cannot easily compact within standard iteration budgets without specialized cross-cluster trajectory bridges.
    \item \textbf{400-Customer Scale (Suboptimal Graceful Degradation)}: Under extreme combinatorial dimensionality ($N=400$), neither heuristic nor hybrid approaches attain BKS. However, Hybrid-DDQN demonstrates superior resistance to search stagnation, securing a vehicle reduction edge of 0.70 to 0.80 vehicles on clustered and random topologies ($c2\_4\_1$ and $r2\_4\_1$), while reducing travel distance by up to 10.66\% ($1309.67\text{ km}$ saved on $r1\_4\_1$).
\end{itemize}

Importantly, our framework is designed as a Quality-Maximizer rather than a search accelerator. While classical ALNS completes 2000 iterations in $1.0\text{s}$ to $15.0\text{s}$ due to raw heuristic evaluation, Hybrid-DDQN incurs an average runtime of $28.9\text{s}$ on Solomon-100, $58.6\text{s}$ on Homberger-200, and $160.9\text{s}$ on Homberger-400. This computational overhead stems from online neural forward/backward passes, prioritized replay buffer management, and periodic exact Set Partitioning solves via HiGHS. On small, lightly constrained instances ($N \le 100$), mature domain-specific genetic algorithms (such as HGS-VRPTW) achieve marginal distance advantages ($+0.13\%$ vs.\ $+1.41\%$) through fast generational crossover; however, Hybrid-DDQN's primary strength lies in maintaining strict operational constraint feasibility, preventing search stagnation, and consistently driving fleet reductions on large, highly constrained topologies ($N=200, 400$).

\subsubsection{Learned Control vs. Handcrafted Rule Thresholds}
\label{sec:disc_rule_vs_learned}
To test whether learned, state-conditioned control is necessary or whether a hand-tuned deterministic policy suffices, Table~\ref{tab:ablation_matrix} (configs 6--8) introduces \textbf{Rule-Macro} (deterministic five-branch tree replacing $\pi_{\text{macro}}$ over the same 12-dimensional state) and \textbf{Rule-Micro} (analogous tree replacing $\pi_{\text{micro}}$ over the 20-dimensional state), each isolated while LAC and HiGHS recombination remain unchanged. On $c2\_2\_1$, Full Hybrid-DDQN converges cleanly to the BKS vehicle floor ($NV=6.00$), whereas heuristic Rule-Macro degrades to $NV=7.20$, Single-Agent RL-LNS degrades to $NV=10.20$, and Rule-Micro severely fails at $NV=12.60$. On $r1\_2\_1$, Full Hybrid-DDQN attains a lower fleet count ($NV=20.40$) than both Rule-Macro ($NV=20.80$) and Rule-Micro ($NV=20.80$). On $R101$, all configurations reliably converge to the minimum fleet size ($NV=19.00$; $TD$ within $0.2\%$ of one another). On the 400-customer scale instance $rc2\_4\_1$, Full Hybrid-DDQN achieves $NV=12.60$, strictly outperforming Rule-Macro ($NV=13.20$), Single-Agent RL-LNS ($NV=13.00$), and Rule-Micro ($NV=13.40$) on the primary fleet minimization metric by eliminating 0.40 to 0.80 vehicles. Taken together, these observations demonstrate that while handcrafted decision trees can perform adequately on simple small-scale instances, state-conditioned Dueling DDQN provides adaptive, general-purpose robustness across large-scale, complex spatiotemporal distributions without requiring manual threshold re-engineering.

\subsection{Managerial Implications for Commercial Logistics}
\label{sec:disc_managerial}

In commercial freight transportation and last-mile parcel distribution, fleet operational expenses are heavily driven by capital expenditures, vehicle licensing, insurance, and driver shift wages rather than marginal fuel consumption \cite{Toth2014, braysy2005vehicle}. Fixed operational overheads typically represent the dominant share of operating budgets, while fuel expenditure contributes a secondary portion.

Consequently, eliminating a single vehicle in a commercial dispatch depot generates substantial recurring cost savings compared to marginal distance reductions. In pre-dispatch planning workflows where routes are generated overnight or during 15- to 30-minute planning windows prior to morning delivery waves, allocating additional computational time to secure vehicle fleet reductions and travel distance savings represents a favorable operational trade-off.

\subsection{Analytical Properties and Complexity Considerations}
\label{sec:disc_theory}

To substantiate the empirical behavior, we examine four analytical properties characterizing the algorithmic mechanics.

\subsubsection{Property 1 (Lexicographic Objective Consistency)}
\label{prop:lex_consistency}
Let $s_1, s_2 \in \mathcal{S}$ be feasible routing solutions. Under the scalar surrogate objective $F(s) = \mathcal{M} \cdot \NV(s) + \TD(s)$, if the penalty weight satisfies $\mathcal{M} > \max_{s \in \mathcal{S}} \TD(s) - \min_{s \in \mathcal{S}} \TD(s)$, then $F(s_1) < F(s_2) \iff s_1 \prec s_2$. In our implementation, $\mathcal{M} = 10^5$ was selected above the maximal total route length across all evaluated benchmark instances, guaranteeing lexicographical equivalence.

\subsubsection{Property 2 (Granular Candidate Space Reduction)}
\label{prop:granular_pruning}
For a VRPTW instance with $N$ customers, dynamic active fleet size $K(s)$, and top-$k$ spatiotemporal candidate filtering, the theoretical pairwise candidate reduction ratio $\Gamma_{\text{prune}}$ for 1-edge insertion moves relative to exhaustive pairwise evaluation is given by:
\begin{equation}
  \Gamma_{\text{prune}}(N, K, k) = 1 - \frac{2 k K(s)}{N(N-1)}
  \label{eq:prop_pruning}
\end{equation}
where $K(s)$ is the number of active routes at state $s$. For $N=200, K=10, k=25$, this yields an evaluation space reduction of $\Gamma_{\text{prune}} = 98.74\%$, and $\Gamma_{\text{prune}} = 99.37\%$ on $N=400, K=20, k=25$. This bound holds strictly for local single-edge substitutions; for multi-customer cluster operations (e.g., Shaw removal), repair complexity is further bounded by candidate adjacency within the surviving route graph.

\subsubsection{Property 3 (Multi-Scale Set Partitioning Penalty Grid)}
\label{prop:penalty_floor}
In Set Partitioning recombination over Route Pool $\mathcal{P}$, the multi-scale penalty grid $\lambda_{\text{pen}} \in \{2.0, 5.0, 12.0, 30.0, 50.0\}\bar{c}$ provides a systematic progression across exploration regimes, ranging from distance preservation to aggressive vehicle elimination, preventing numerical stalling during exact MIP solves.

\subsubsection{Property 4 (Downstream Lookahead under Learned Acceptance)}
\label{prop:lac_behavior}
In classical Simulated Annealing, the acceptance probability for a non-improving move $\Delta C > 0$ decays exponentially as temperature cools ($P_{\text{SA}} = \exp(-\Delta C / T) \to 0$ as $T \to 0$), restricting late-stage search to strictly downhill transitions. In contrast, the Learned Acceptance Criterion augments Simulated Annealing by evaluating structural transition promise based on delayed lookahead labeling ($H_{\text{lac}} = 50$). When the classifier predicts a high probability of future incumbent discovery ($p_t \ge \tau_{\text{acc}}$), non-improving intermediate moves are retained, facilitating plateau escape during deep cooling stages.

\subsection{Limitations and Threats to Validity}
\label{sec:disc_limitations}

While Tri-Level Hybrid DDQN-ALNS exhibits strong empirical performance, five primary limitations should be noted:
\begin{enumerate}
    \item \textbf{Computational Latency on Small Instances}: On small-scale problems ($N \le 50$), classical heuristics find near-optimal solutions in fractions of a second. The online neural training and replay buffer maintenance of Hybrid-DDQN introduce unnecessary overhead in this regime without substantial solution quality gains.
    \item \textbf{MIP Solver Dependency}: The performance of the matheuristic recombination tier depends on the availability and efficiency of an underlying MIP solver (HiGHS). On dense route pools ($|\mathcal{P}| > 2000$), solve time limits ($4.0\text{s}$) must be enforced to prevent solver stalling.
    \item \textbf{Late-Stage Lookahead Horizon Truncation in LAC}: The LAC's 50-iteration lookahead horizon is calibrated for local plateau escapes and becomes increasingly conservative in late-stage search as class imbalance ($y_t = 1$ rate $< 2\%$) biases the online MLP toward low acceptance probabilities. Compound fleet-reduction restructuring requiring $> 50$ iterations is explicitly delegated to the Macro Controller's HiGHS and thermal reheating modes.
    \item \textbf{Fixed-Iteration vs. Equal Wall-Clock Comparison}: Our primary benchmark suite evaluated algorithms under identical iteration limits ($T_{\max} = 2000$). Because learning-augmented components introduce per-step computational overhead, Hybrid-DDQN utilizes a larger wall-clock budget than baseline ALNS.
    \item \textbf{Online Adaptation Scope}: The framework relies on online instance-adaptive reinforcement learning during the search trajectory. For sub-second dynamic dispatch scenarios with extreme real-time network distortions, offline pretrained policy initialization may be required.
\end{enumerate}
